\section{Canonical ratio sections and the original group comparison}

For a scheme $N$, retain actual invertible ideal modules $I,J$ and their
original inclusions into $\mathcal O_N$. A canonical ratio section is
the inverse numerator inclusion tensored with the evaluation transpose
of the denominator inclusion. It is defined wherever the numerator
inclusion is invertible. This construction concerns actual module
sections before taking any valuation or norm.

\begin{lemma}[Pulled evaluation transposes preserve the actual functional]
\label{lem:original-pl-ratio-pulled-functional-transport}
\lean{EnvelopeProperty.PLOriginalModelCompatibility.RatioSectionTransport.pulledDualTranspose_pairing, EnvelopeProperty.PLOriginalModelCompatibility.RatioSectionTransport.evaluationDualTransport_dualTranspose}
\leanok
Let $f,g:T\to N$ and let $e:f^*J\simeq g^*J$ preserve the original
pulled inclusions into $\mathcal O_T$. The canonical evaluation-preserving
transport $f^*(J^\vee)\simeq g^*(J^\vee)$ sends the pulled transpose
of the same original denominator inclusion to its other pullback.
\end{lemma}
\begin{proof}
The defining evaluation identity for the original transpose remains an
identity after pullback. The strong monoidal tensorator is natural,
and its unit comparison is the canonical pullback of the structure
sheaf. The monoidal right-unit identity therefore identifies this
pulled identity with the evaluation identity of the pulled functional.

Tensor the proposed transported section by the invertible module
$g^*J$ and evaluate. The defining identity for evaluation-preserving
dual transport changes this into evaluation against $f^*J$, preceded
by $e^{-1}$. The original transpose identity makes it the functional
$e^{-1}\circ f^*\iota_J$. Preservation of the original inclusion makes
this exactly $g^*\iota_J$. The other pulled transpose has the same
evaluation. Evaluation is an isomorphism, and tensoring by an invertible
module is faithful. Cancellation proves equality of the actual sections.
\end{proof}

\begin{theorem}[Actual canonical ratio sections are preserved]
\label{thm:original-pl-canonical-ratio-section-transport}
\lean{EnvelopeProperty.PLOriginalModelCompatibility.RatioSectionTransport.inverseInclusion_transport, EnvelopeProperty.PLOriginalModelCompatibility.RatioSectionTransport.canonicalPulledRatioSection, EnvelopeProperty.PLOriginalModelCompatibility.RatioSectionTransport.evaluationRatioComparison_canonicalSection}
\leanok
\uses{lem:original-pl-ratio-pulled-functional-transport}
Let $f,g:T\to N$ have invertible pulled numerator inclusions, and
let $e_I:f^*I\simeq g^*I$, $e_J:f^*J\simeq g^*J$ preserve the two
original pulled inclusions. The canonical tensor-dual ratio comparison
constructed from $e_I,e_J$ sends the actual canonical ratio section
for $f$ to the actual canonical ratio section for $g$.
\end{theorem}
\begin{proof}
Compose the transported inverse numerator inclusion with the original
inclusion of $g^*I$. Inclusion preservation makes this composite the
identity, so cancellation by that invertible inclusion identifies it
with the inverse inclusion of $g^*I$. The preceding lemma identifies
the transported denominator transpose. Tensor these two section
identities using the original structure-sheaf multiplication. Naturality
of the canonical coherent-to-explicit tensor comparison identifies
this tensor transport with the defining whole ratio comparison. This
proves equality of the sections of the unchanged pulled ratio module.
\end{proof}

\begin{corollary}[The canonical group ratio comparison preserves its section]
\label{cor:original-pl-group-canonical-ratio-section}
\lean{EnvelopeProperty.PLOriginalModelCompatibility.RatioSectionTransport.pullbackIsoAlong_inclusion, EnvelopeProperty.PLOriginalModelCompatibility.RatioSectionTransport.cartierRatioGroupLinearization_canonicalSection, EnvelopeProperty.PLOriginalModelCompatibility.RatioSectionTransport.offSupportInclusion_isIso, EnvelopeProperty.PLOriginalModelCompatibility.RatioSectionTransport.canonicalPulledRatioSection_eq_localSection, EnvelopeProperty.PLOriginalModelCompatibility.RatioSectionTransport.offSupportFactorInclusion_isIso, EnvelopeProperty.PLOriginalModelCompatibility.RatioSectionTransport.canonicalPulledRatioSection_eq_sectionViaFactor_localSection, EnvelopeProperty.PLOriginalModelCompatibility.RatioSectionTransport.cartierRatioGroupLinearization_localSection}
\leanok
\uses{thm:original-pl-canonical-ratio-section-transport}
For the actual action and projection of a group acting over a field,
retain the same invariant Cartier ideals and their canonical
inclusion-preserving group comparisons. After any further base change
on which the pulled numerator inclusions are invertible, the canonical
ratio group comparison preserves the actual canonical ratio section.
No section-matching or norm-matching identity is an input.
\end{corollary}
\begin{proof}
The original canonical factor comparisons preserve both entire ideal
inclusions. Pulling back these identities and using the canonical
composition comparison for the structure sheaf proves the same
inclusion identities after the further base change. Canonical ratio
comparisons commute with this arbitrary base change. Apply the preceding
theorem to these pulled comparisons. On an original off-support open
the inverse pulled inclusion is its original canonical inclusion frame,
so these sections are exactly the original local Cartier ratio sections.
The existing residue-fibre identification and model-weight identities
therefore apply to these same sections, without replacing them by
sections with merely equal norms.
\end{proof}
